%!TeX root=../persuasiontop.tex
\chapter[Chapter \thechapter]{}

\begin{figure}[t!]
\centering
\includegraphics[width=\linewidth]{nouveau07}
\end{figure}

\lettrine[lines=4,lraise=0.3]{A}{} very few days more, and Captain Wentworth was known to be at Kellynch, and Mr Musgrove had called on him, and come back warm in his praise, and he was engaged with the Crofts to dine at Uppercross, by the end of another week. It had been a great disappointment to Mr Musgrove to find that no earlier day could be fixed, so impatient was he to shew his gratitude, by seeing Captain Wentworth under his own roof, and welcoming him to all that was strongest and best in his cellars. But a week must pass; only a week, in Anne's reckoning, and then, she supposed, they must meet; and soon she began to wish that she could feel secure even for a week.

Captain Wentworth made a very early return to Mr Musgrove's civility, and she was all but calling there in the same half hour. She and Mary were actually setting forward for the Great House, where, as she afterwards learnt, they must inevitably have found him, when they were stopped by the eldest boy's being at that moment brought home in consequence of a bad fall. The child's situation put the visit entirely aside; but she could not hear of her escape with indifference, even in the midst of the serious anxiety which they afterwards felt on his account.

His collar-bone was found to be dislocated, and such injury received in the back, as roused the most alarming ideas. It was an afternoon of distress, and Anne had every thing to do at once; the apothecary to send for, the father to have pursued and informed, the mother to support and keep from hysterics, the servants to control, the youngest child to banish, and the poor suffering one to attend and soothe; besides sending, as soon as she recollected it, proper notice to the other house, which brought her an accession rather of frightened, enquiring companions, than of very useful assistants.

Her brother's return was the first comfort; he could take best care of his wife; and the second blessing was the arrival of the apothecary. Till he came and had examined the child, their apprehensions were the worse for being vague; they suspected great injury, but knew not where; but now the collar-bone was soon replaced, and though Mr Robinson felt and felt, and rubbed, and looked grave, and spoke low words both to the father and the aunt, still they were all to hope the best, and to be able to part and eat their dinner in tolerable ease of mind; and then it was, just before they parted, that the two young aunts were able so far to digress from their nephew's state, as to give the information of Captain Wentworth's visit; staying five minutes behind their father and mother, to endeavour to express how perfectly delighted they were with him, how much handsomer, how infinitely more agreeable they thought him than any individual among their male acquaintance, who had been at all a favourite before. How glad they had been to hear papa invite him to stay dinner, how sorry when he said it was quite out of his power, and how glad again when he had promised in reply to papa and mamma's farther pressing invitations to come and dine with them on the morrow—actually on the morrow; and he had promised it in so pleasant a manner, as if he felt all the motive of their attention just as he ought. And in short, he had looked and said everything with such exquisite grace, that they could assure them all, their heads were both turned by him; and off they ran, quite as full of glee as of love, and apparently more full of Captain Wentworth than of little Charles.

The same story and the same raptures were repeated, when the two girls came with their father, through the gloom of the evening, to make enquiries; and Mr Musgrove, no longer under the first uneasiness about his heir, could add his confirmation and praise, and hope there would be now no occasion for putting Captain Wentworth off, and only be sorry to think that the cottage party, probably, would not like to leave the little boy, to give him the meeting. »Oh no; as to leaving the little boy,« both father and mother were in much too strong and recent alarm to bear the thought; and Anne, in the joy of the escape, could not help adding her warm protestations to theirs.

Charles Musgrove, indeed, afterwards, shewed more of inclination; »the child was going on so well, and he wished so much to be introduced to Captain Wentworth, that, perhaps, he might join them in the evening; he would not dine from home, but he might walk in for half an hour.« But in this he was eagerly opposed by his wife, with »Oh! no, indeed, Charles, I cannot bear to have you go away. Only think if anything should happen?«

The child had a good night, and was going on well the next day. It must be a work of time to ascertain that no injury had been done to the spine; but Mr Robinson found nothing to increase alarm, and Charles Musgrove began, consequently, to feel no necessity for longer confinement. The child was to be kept in bed and amused as quietly as possible; but what was there for a father to do? This was quite a female case, and it would be highly absurd in him, who could be of no use at home, to shut himself up. His father very much wished him to meet Captain Wentworth, and there being no sufficient reason against it, he ought to go; and it ended in his making a bold, public declaration, when he came in from shooting, of his meaning to dress directly, and dine at the other house.

»Nothing can be going on better than the child,« said he; »so I told my father, just now, that I would come, and he thought me quite right. Your sister being with you, my love, I have no scruple at all. You would not like to leave him yourself, but you see I can be of no use. Anne will send for me if anything is the matter.«

Husbands and wives generally understand when opposition will be vain. Mary knew, from Charles's manner of speaking, that he was quite determined on going, and that it would be of no use to teaze him. She said nothing, therefore, till he was out of the room, but as soon as there was only Anne to hear—

»So you and I are to be left to shift by ourselves, with this poor sick child; and not a creature coming near us all the evening! I knew how it would be. This is always my luck. If there is anything disagreeable going on men are always sure to get out of it, and Charles is as bad as any of them. Very unfeeling! I must say it is very unfeeling of him to be running away from his poor little boy. Talks of his being going on so well! How does he know that he is going on well, or that there may not be a sudden change half an hour hence? I did not think Charles would have been so unfeeling. So here he is to go away and enjoy himself, and because I am the poor mother, I am not to be allowed to stir; and yet, I am sure, I am more unfit than anybody else to be about the child. My being the mother is the very reason why my feelings should not be tried. I am not at all equal to it. You saw how hysterical I was yesterday.«

»But that was only the effect of the suddenness of your alarm—of the shock. You will not be hysterical again. I dare say we shall have nothing to distress us. I perfectly understand Mr Robinson's directions, and have no fears; and indeed, Mary, I cannot wonder at your husband. Nursing does not belong to a man; it is not his province. A sick child is always the mother's property: her own feelings generally make it so.«

»I hope I am as fond of my child as any mother, but I do not know that I am of any more use in the sick-room than Charles, for I cannot be always scolding and teazing the poor child when it is ill; and you saw, this morning, that if I told him to keep quiet, he was sure to begin kicking about. I have not nerves for the sort of thing.«

»But, could you be comfortable yourself, to be spending the whole evening away from the poor boy?«

»Yes; you see his papa can, and why should not I? Jemima is so careful; and she could send us word every hour how he was. I really think Charles might as well have told his father we would all come. I am not more alarmed about little Charles now than he is. I was dreadfully alarmed yesterday, but the case is very different to-day.«

»Well, if you do not think it too late to give notice for yourself, suppose you were to go, as well as your husband. Leave little Charles to my care. Mr and Mrs Musgrove cannot think it wrong while I remain with him.«

»Are you serious?« cried Mary, her eyes brightening. »Dear me! that's a very good thought, very good, indeed. To be sure, I may just as well go as not, for I am of no use at home—am I? and it only harasses me. You, who have not a mother's feelings, are a great deal the properest person. You can make little Charles do anything; he always minds you at a word. It will be a great deal better than leaving him only with Jemima. Oh! I shall certainly go; I am sure I ought if I can, quite as much as Charles, for they want me excessively to be acquainted with Captain Wentworth, and I know you do not mind being left alone. An excellent thought of yours, indeed, Anne. I will go and tell Charles, and get ready directly. You can send for us, you know, at a moment's notice, if anything is the matter; but I dare say there will be nothing to alarm you. I should not go, you may be sure, if I did not feel quite at ease about my dear child.«

The next moment she was tapping at her husband's dressing-room door, and as Anne followed her up stairs, she was in time for the whole conversation, which began with Mary's saying, in a tone of great exultation—

»I mean to go with you, Charles, for I am of no more use at home than you are. If I were to shut myself up for ever with the child, I should not be able to persuade him to do anything he did not like. Anne will stay; Anne undertakes to stay at home and take care of him. It is Anne's own proposal, and so I shall go with you, which will be a great deal better, for I have not dined at the other house since Tuesday.«

»This is very kind of Anne,« was her husband's answer, »and I should be very glad to have you go; but it seems rather hard that she should be left at home by herself, to nurse our sick child.«

Anne was now at hand to take up her own cause, and the sincerity of her manner being soon sufficient to convince him, where conviction was at least very agreeable, he had no farther scruples as to her being left to dine alone, though he still wanted her to join them in the evening, when the child might be at rest for the night, and kindly urged her to let him come and fetch her, but she was quite unpersuadable; and this being the case, she had ere long the pleasure of seeing them set off together in high spirits. They were gone, she hoped, to be happy, however oddly constructed such happiness might seem; as for herself, she was left with as many sensations of comfort, as were, perhaps, ever likely to be hers. She knew herself to be of the first utility to the child; and what was it to her if Frederick Wentworth were only half a mile distant, making himself agreeable to others?

She would have liked to know how he felt as to a meeting. Perhaps indifferent, if indifference could exist under such circumstances. He must be either indifferent or unwilling. Had he wished ever to see her again, he need not have waited till this time; he would have done what she could not but believe that in his place she should have done long ago, when events had been early giving him the independence which alone had been wanting.

Her brother and sister came back delighted with their new acquaintance, and their visit in general. There had been music, singing, talking, laughing, all that was most agreeable; charming manners in Captain Wentworth, no shyness or reserve; they seemed all to know each other perfectly, and he was coming the very next morning to shoot with Charles. He was to come to breakfast, but not at the Cottage, though that had been proposed at first; but then he had been pressed to come to the Great House instead, and he seemed afraid of being in Mrs Charles Musgrove's way, on account of the child, and therefore, somehow, they hardly knew how, it ended in Charles's being to meet him to breakfast at his father's.

Anne understood it. He wished to avoid seeing her. He had inquired after her, she found, slightly, as might suit a former slight acquaintance, seeming to acknowledge such as she had acknowledged, actuated, perhaps, by the same view of escaping introduction when they were to meet.

The morning hours of the Cottage were always later than those of the other house, and on the morrow the difference was so great that Mary and Anne were not more than beginning breakfast when Charles came in to say that they were just setting off, that he was come for his dogs, that his sisters were following with Captain Wentworth; his sisters meaning to visit Mary and the child, and Captain Wentworth proposing also to wait on her for a few minutes if not inconvenient; and though Charles had answered for the child's being in no such state as could make it inconvenient, Captain Wentworth would not be satisfied without his running on to give notice.

Mary, very much gratified by this attention, was delighted to receive him, while a thousand feelings rushed on Anne, of which this was the most consoling, that it would soon be over. And it was soon over. In two minutes after Charles's preparation, the others appeared; they were in the drawing-room. Her eye half met Captain Wentworth's, a bow, a curtsey passed; she heard his voice; he talked to Mary, said all that was right, said something to the Miss Musgroves, enough to mark an easy footing; the room seemed full, full of persons and voices, but a few minutes ended it. Charles shewed himself at the window, all was ready, their visitor had bowed and was gone, the Miss Musgroves were gone too, suddenly resolving to walk to the end of the village with the sportsmen: the room was cleared, and Anne might finish her breakfast as she could.

»It is over! it is over!« she repeated to herself again and again, in nervous gratitude. »The worst is over!«

Mary talked, but she could not attend. She had seen him. They had met. They had been once more in the same room.

Soon, however, she began to reason with herself, and try to be feeling less. Eight years, almost eight years had passed, since all had been given up. How absurd to be resuming the agitation which such an interval had banished into distance and indistinctness! What might not eight years do? Events of every description, changes, alienations, removals—all, all must be comprised in it, and oblivion of the past— how natural, how certain too! It included nearly a third part of her own life.

Alas! with all her reasoning, she found, that to retentive feelings eight years may be little more than nothing.

Now, how were his sentiments to be read? Was this like wishing to avoid her? And the next moment she was hating herself for the folly which asked the question.

On one other question which perhaps her utmost wisdom might not have prevented, she was soon spared all suspense; for, after the Miss Musgroves had returned and finished their visit at the Cottage she had this spontaneous information from Mary:—

»Captain Wentworth is not very gallant by you, Anne, though he was so attentive to me. Henrietta asked him what he thought of you, when they went away, and he said, »You were so altered he should not have known you again.««

Mary had no feelings to make her respect her sister's in a common way, but she was perfectly unsuspicious of being inflicting any peculiar wound.

»Altered beyond his knowledge.« Anne fully submitted, in silent, deep mortification. Doubtless it was so, and she could take no revenge, for he was not altered, or not for the worse. She had already acknowledged it to herself, and she could not think differently, let him think of her as he would. No: the years which had destroyed her youth and bloom had only given him a more glowing, manly, open look, in no respect lessening his personal advantages. She had seen the same Frederick Wentworth.

»So altered that he should not have known her again!« These were words which could not but dwell with her. Yet she soon began to rejoice that she had heard them. They were of sobering tendency; they allayed agitation; they composed, and consequently must make her happier.

Frederick Wentworth had used such words, or something like them, but without an idea that they would be carried round to her. He had thought her wretchedly altered, and in the first moment of appeal, had spoken as he felt. He had not forgiven Anne Elliot. She had used him ill, deserted and disappointed him; and worse, she had shewn a feebleness of character in doing so, which his own decided, confident temper could not endure. She had given him up to oblige others. It had been the effect of over-persuasion. It had been weakness and timidity.

He had been most warmly attached to her, and had never seen a woman since whom he thought her equal; but, except from some natural sensation of curiosity, he had no desire of meeting her again. Her power with him was gone for ever.

It was now his object to marry. He was rich, and being turned on shore, fully intended to settle as soon as he could be properly tempted; actually looking round, ready to fall in love with all the speed which a clear head and a quick taste could allow. He had a heart for either of the Miss Musgroves, if they could catch it; a heart, in short, for any pleasing young woman who came in his way, excepting Anne Elliot. This was his only secret exception, when he said to his sister, in answer to her suppositions:—

»Yes, here I am, Sophia, quite ready to make a foolish match. Anybody between fifteen and thirty may have me for asking. A little beauty, and a few smiles, and a few compliments to the navy, and I am a lost man. Should not this be enough for a sailor, who has had no society among women to make him nice?«

He said it, she knew, to be contradicted. His bright proud eye spoke the conviction that he was nice; and Anne Elliot was not out of his thoughts, when he more seriously described the woman he should wish to meet with. »A strong mind, with sweetness of manner,« made the first and the last of the description.

»That is the woman I want,« said he. »Something a little inferior I shall of course put up with, but it must not be much. If I am a fool, I shall be a fool indeed, for I have thought on the subject more than most men.«